\section{Limitations and threats to validity}
\label{sec:limitations}

Every figure in \cref{sec:results} is checked against exact roots by a gate that
fails on a missed collision or a late time of impact, and the data behind each
table is published alongside the code. Three kinds of limitation nevertheless
lie outside what those numbers establish: the reach of the ground truth they are
measured against, the reach of the measurement itself, and the parts of the
pipeline that carry no measurement.

\subsection{What the ground truth reaches}

The dataset ships symbolic roots for its curated \emph{query sets} and not for
the mesh frames, so everything in \cref{sec:results:cons,sec:results:ref} is
measured on the query sets. The mesh path is checked only for
\emph{agreement}: the answer computed over the mesh is compared against the
earliest exact root of the corresponding query set, and a disagreement would
mean the two paths were handed different geometry. This is a tripwire on the
inputs, not a verification of the kernel.

The two paths do differ in the geometry they carry. The mesh container stores
coordinates in single precision, and armadillo-rollers needs $52$ significand
bits to represent its dyadic rationals exactly. The resulting divergence between
mesh and query geometry is two-sided, and it ranges from $5\times10^{-8}$ to
$4.6\times10^{-5}$. An accuracy claim tighter than that, made on the mesh path,
measures the rounding and not the search.

Deliberately inflicting that rounding shows the size of its effect. On the same
$98{,}757$ armadillo-rollers edge-edge queries, \textsc{Tight} reports no missed
collision and no late time of impact on the shipped double-precision geometry.
Once the input is first rounded through single precision, it reports $4$ missed
and $25{,}485$ late. A control settles where the fault lies: a host kernel
bit-identical to TightInclusion, already computing in double, reports zero late
on the shipped geometry and $4{,}543$ late under the same treatment. The cause
is the input rounding and not either kernel. Single-precision-geometry runs are
therefore reported separately and never gate the build, and the comparison is
the empirical form of the argument in \cref{sec:parallel:precision}.
\textsc{Tight} is hurt about fourteen times more than \textsc{Relaxed} here
\emph{because} it reports closer to the true root: a perturbed root crosses a
tight answer far more often than it crosses a loose one.

\subsection{What a measurement establishes}

When the caller asks only for the earliest time of impact, workers prune against
a shared atomic minimum, so the exact value depends on the order in which they
publish. Every value observed is at or before the true root, and the guarantee
is unaffected. The last digits do move between runs, however, so a test that
pins a time of impact to better than about $10^{-6}$ will be flaky. The same
applies to the quad path.

Non-determinism of that kind can also hide a defect. During this evaluation one
repeat of one scene showed the device \textsc{Relaxed} kernel missing $2$
collisions and reporting $5$ late of $15{,}957$ queries. The next repeat over
identical input was clean, as were the other five scenes and the other three
modes across the remaining $65{,}079{,}040$ checks of that run. The cause was a
single index, which pruned the root box against the first query's answer in
place of its own. This occurs only when the first query reports exactly zero, so
whether a collision was missed depended on block scheduling. Twelve repeats of
the case failed five times before the fix and none after it, and every device
figure in this paper is taken after the fix.

The episode bounds what a check can establish, not what the kernel does. A
defect surfacing in about two runs in five is invisible to a spot check and
survives a benchmark that reports only agreement; separating it from noise takes
a signed comparison against exact roots, run over every query and repeated.

A last caveat concerns the timings, not the answers. The driver does more per
case than the pipeline does: it loads the frame, classifies every candidate pair
against the expected set, and compares every answer against an exact root. A
case costs between $0.7$ and $3.3\,\mathrm{s}$ of wall clock across the six
scenes, against measured work of a few tens of milliseconds. The timers exclude
all of it, so the figures are sound as timings, although the harness remains no
model of what the pipeline costs inside a solver.

\subsection{What was not measured}

Three parts of the pipeline carry no measurement here, and each bounds a claim
made earlier.

Ordering each cell's entries on the axis the grid does not use is implemented
on the host only, where it is measured and does not pay
(\cref{sec:bp:self}); the device therefore carries no such kernel and the
question is open on that processor.

The quad path is described in \cref{sec:np:quad} and evaluated nowhere. Every
scene in this benchmark is a triangle mesh, and the error constant of
\cref{sec:prelim:err} for the bilinear form is uncertified. What is established
for that path is correctness against a densely sampled reference root, and
nothing about its speed or its accuracy at scale.

Most consequentially, the closest prior work is cited and not measured against.
\citet{belgrod2025toi} evaluate a GPU CCD pipeline on these six scenes against
this ground truth, and theirs, not TightInclusion, is the state of the art in
\emph{parallel} CCD. TightInclusion answers the question this paper puts first,
which is whether the answer is right and how close it falls to the exact root,
and it answers that as a certified reference running unmodified on the same
hardware. The remaining question, whether this parallel design is faster than
that parallel design, is left open. Settling it requires building that pipeline
on \textsc{GH200}, matching the two on tolerance, depth cap and output mode so
the searches do comparable work, and timing both inside one allocation. Until
then the speedups of \cref{sec:results:ref} are margins over a host reference
and not over prior parallel work.

The counter measurements of \cref{sec:parallel:limits} are narrower than the
timings they explain. They cover all six scenes, but only one narrow-phase mode
and only the earliest-impact kernel. The device figures are taken under a
profiler that serialises launches, so they support statements about ratios and
none about absolute time. The Grace rows are whole-process counters and not
per-phase ones. The pipeline row comes from a build with the device path
compiled out, so it is host work throughout. The narrow-phase row comes from the
oracle, which also drives the device modes, so its instruction rate is a floor
and not an estimate. These counters establish which resource each processor is
short of, a coarser claim than a per-phase model would make.

Beyond the pipeline itself, three things are absent. There is no
minimum-separation CCD, which solvers in the barrier family generally want, and
no multi-GPU execution. Nor is there any measurement inside a solver, so the
end-to-end benefit of the speedups reported here is inferred and not
established.
